Work in $\Cl(2,4)$ with the generators of Theorem \claimref{synthesis.gtg-tractor-gauge-field-reproduces-tractor-connection} ($e_{0}$, $e_{+}$ squaring to $+1$; $e_{1}, e_{2}, e_{3}, e_{-}$ to $-1$; $\nO = \tfrac12(e_{+} + e_{-})$, $\nInf = e_{+} - e_{-}$), so $\nu = e_{+}$ and $Z = e_{0}e_{+}$ in the notation of Proposition \claimref{synthesis.attachment-conformal-energy-is-generated-by-z}. By Proposition \claimref{field-theory.where-each-lives}, $\CS$ is the minimal left ideal $\Cl^{+}(2,4)P$, with complex unit $I_{6} = I_{4}E$ ($E = e_{+}e_{-}$) and the splitting $\CS = L_{0} \oplus L_{\infty}$ into the two chiral doublets, on which $I_{6} = \pm I_{4}$. So $\Cl^{+}(2,4)$, a complex algebra of complex dimension $16$, acts on itself from the left as four copies of $\CS$, and the subspace of $\CS$ annihilated by a set of bivectors has complex dimension one eighth of the real dimension of their common left kernel on $\Cl^{+}(2,4)$; the checks compute that kernel.

\emph{Transport and integrability.} The tractor gauge field acts on a $\CS$-valued field through the spin action, $D^{T}_{a}\psi = a\cdot\nabla\psi + \tfrac12\Omega_{T}(a)\psi$, whose commutator on the plane of $a, b$ is $\tfrac12 F_{T}(a\wedge b)\psi$ with $F_{T}$ the field strength of $\Omega_{T}$ (the standard curvature of a gauge field with the gauge law of $\Omega$, Fact \claimref{synthesis.imported-gtg-dictionary}). A parallel field therefore satisfies $F_{T}(a\wedge b)\psi = 0$ at every point and for every plane, and a parallel adjoint tractor $X$ (a bivector field with $D^{T}_{a}X = a\cdot\nabla X + \Omega_{T}(a)\times X = 0$) satisfies $F_{T}(a\wedge b)\times X = 0$. Conversely, where $F_{T} = 0$ on a simply connected region the connection is flat and the parallel sections of $\CS$ form a space of complex dimension $4$, on which the constant $\mathfrak{u}(2,2)$ action, the form $h$ and $Z$ of the attachment are carried unchanged.

\emph{(i) Ricci blindness.} By Theorem \claimref{synthesis.gtg-tractor-gauge-field-reproduces-tractor-connection}(iv), on a torsion-free field the part of $F_{T}$ with no derivative of the frame or of $\Sch$ is $-\mathcal{R}(A\wedge B) + \Theta(A)\times\Theta(B)$, whose spacetime block is $-\mathcal{W}(A\wedge B)$ with no other block; the derivative part is the Cotton term in the $\nO$ block. The check recomputes this with $\mathcal{R}(A\wedge B) = \mathcal{W}(A\wedge B) + \Sch(A)\wedge B + A\wedge\Sch(B)$ for six Weyl maps (one of each Petrov type) and two non-Einstein symmetric $\Sch$: the Schouten map cancels on all six coordinate planes. Each Weyl map is built as an $I_{4}$-linear map from a complex symmetric traceless $3\times3$ matrix $M$ in the basis $\sigma_{k} = e_{k}e_{0}$; the check confirms pair symmetry, vanishing Ricci trace, the first Bianchi identity, and the type from the minimal polynomial of $M$ (O: $M = 0$; N: $M^{2} = 0$; III: $M^{3} = 0 \neq M^{2}$; otherwise a repeated eigenvalue exactly when $(\mathrm{tr}M^{2})^{3} = 6(\mathrm{tr}M^{3})^{2}$, D or II as $(M - \lambda)(M + 2\lambda)$ vanishes or not, I if no repeat).

\emph{(ii) The annihilated subspace.} (Two-spinor reading, labelled: on $L_{0} \oplus L_{\infty}$ a conformal spinor is a pair $(\omega^{A}, \pi_{A'})$, and the Weyl bivectors act on $\omega$ through the Weyl spinor $\Psi_{ABCD}$ and on $\pi$ through its conjugate.) The common annihilator of the Weyl bivectors in $\CS$ is therefore $\{\omega : \Psi_{ABCD}\omega^{D} = 0\} \oplus \{\pi : \bar\Psi_{A'B'C'D'}\pi^{D'} = 0\}$. If $\omega \neq 0$ and $\Psi_{ABCD}\omega^{D} = 0$, complete $o = \omega$ to a spin frame $(o, \iota)$: every component of $\Psi$ with at least one contraction with $o$ vanishes, so $\Psi_{ABCD} = \Psi_{4}\,o_{A}o_{B}o_{C}o_{D}$, type N with principal spinor $o$ (or zero); conversely at a type N point $\Psi_{ABCD}\omega^{D} = 0$ exactly for $\omega$ proportional to $o$. This is the standard integrability condition of the twistor equation, which restricts the field to Petrov types N and O (Fact \claimref{synthesis.imported-twistor-equation-integrability-petrov-n-or-o}). So the complex dimension is $4$ at type O, $2$ at type N and $0$ at types I, II, D, III. The check computes the common left kernel of the six field strengths $F_{T}(e_{m}\wedge e_{n})$ for each Petrov sample: real dimension $32, 16, 0, 0, 0, 0$, that is complex dimension $4, 2, 0, 0, 0, 0$ on $\CS$; for type N it coincides with the left kernel of the repeated principal null vector $k = e_{0} + e_{3}$ (which every type N curvature bivector annihilates under the commutator product) and meets the left kernel of $m = e_{0} - e_{3}$ only in zero. The samples are general: type N and type III normal forms are unique up to a Lorentz rotor and a complex scale, neither of which changes the kernel dimension, and the spinor argument above covers types I, II and D for every value of their invariants. A Cotton block $\nO\wedge c$ adds the condition of being annihilated by it, which can only cut the subspace; in the samples it keeps the type N subspace for $c$ along $k$ and cuts it to complex dimension $1$ for $c = e_{1}, e_{0}, m$. Where $\mathcal{W}$ vanishes on a region the Cotton term vanishes too and $F_{T} = 0$: this converse is proved on twistor-flat fields in \claimref{synthesis.twistor-flat-fields-give-nordstrom-not-einstein}, clause (i), from Weyl's theorem (stored as an imported Fact, correspondence.imported-weyls-theorem) and the conformal invariance of the tractor connection, and is not used in the checks. On every Einstein space ($\Sch = \tfrac{\Lambda}{6}\id$, parallel) the Cotton term vanishes; under Einstein's equations with matter it is in general nonzero.

\emph{(iii) No conformal-energy generator.} At points of type II and D, and of type I with nonzero eigenvalues, the matrix $M$ is invertible, so the Weyl bivectors span all of $\so(1,3)$, and the bivectors of $\Cl(2,4)$ commuting with $\so(1,3)$ are the multiples of $E = e_{+}e_{-}$. Type III, and type I with a zero eigenvalue, have normal forms unique up to a Lorentz rotor and a complex scale; for both the check finds the same one-dimensional commutant, spanned by $E$. At type N the commutant is the five-dimensional span of $k\wedge e_{1}$, $k\wedge e_{2}$, $k\wedge e_{+}$, $k\wedge e_{-}$ and $E$; at type O it is all fifteen bivectors. In each case with $\mathcal{W} \neq 0$ every element of the commutant annihilates the null vector $k$ under the commutator product ($E$ annihilates every spacetime vector). The generator $Z = e_{0}e_{+}$ has $Z^{2} = -1$ and annihilates exactly the span of $e_{1}, e_{2}, e_{3}, e_{-}$, a negative-definite $4$-plane; a rotor $g$ is an isometry, so $\lambda\,gZ\tilde g$ ($\lambda \neq 0$) annihilates exactly the negative-definite plane $g(\mathrm{span}\{e_{1}, e_{2}, e_{3}, e_{-}\})$, which contains no null vector. Hence no nonzero multiple of a conjugate of $Z$ lies in the commutant at a point with $\mathcal{W} \neq 0$, and no parallel adjoint tractor through such a point is a conformal-energy generator; at type O, $Z$ commutes with the curvature (computed).

\emph{Reading.} The attachment's one-particle structure is $\CS$ with its invariant form $h$, the $\mathfrak{u}(2,2)$ action, and the grading $Z$ that turns $h$ into the Hilbert form and generates the positive conformal energy of the massless, conformally invariant modules ($Z$ is not a symmetry of a massive field even on a flat region) (Propositions \claimref{field-theory.twistor-attachment}, \claimref{synthesis.attachment-conformal-energy-is-generated-by-z}). By (ii) and (iii) the tractor transport carries it across a region only where $\mathcal{W} = 0$; where $\mathcal{W} \neq 0$ at most a two-dimensional subspace of $\CS$ is parallel, and no conformal-energy generator is. So on a general gravitational field the attachment's tractor transport does not construct the one-particle space and supplies no state: a choice of positive frequency has to come from another construction, which this result does not exclude. The companion result on the Hadamard class and null planes (\claimref{synthesis.bulk-hadamard-states-restrict-to-the-generator-condition-on-flat-null-planes}) is proved on the flat field only and says nothing about fields with $\mathcal{W} \neq 0$. That the field algebra itself, the CAR algebra over the Dirac solution space, exists on every field with a Cauchy surface is standard (Fact \claimref{synthesis.imported-dimock-dirac-car-algebra-on-a-gravitational-field}) and is not used here; what is missing off Weyl-flat regions is a state on that algebra, not the algebra. Transport-convenient: the result is the standard integrability condition of the twistor equation, reproduced in the series' tractor gauge field and read through the attachment's $Z$; the series' transport adds the reading, not the mathematics.
